%!TEX root = ../template.tex
%%%%%%%%%%%%%%%%%%%%%%%%%%%%%%%%%%%%%%%%%%%%%%%%%%%%%%%%%%%%%%%%%%%%
%% abstract-en.tex
%% NOVA thesis document file
%%
%% Abstract in English
%%%%%%%%%%%%%%%%%%%%%%%%%%%%%%%%%%%%%%%%%%%%%%%%%%%%%%%%%%%%%%%%%%%%

\typeout{NT FILE abstract-en.tex}%

Standard sports wearables rely predominantly on kinematic data from Inertial Measurement Units (\glspl{IMU}), missing the physiological context of the muscles to classify complex, high-intensity athletic movements. 
Addressing this limitation, this dissertation presents the design, implementation, and evaluation of a multi-modal wearable ecosystem tailored for basketball analytics. The custom hardware architecture fuses 9-axis 
\gls{IMU} kinematics with Electromyography (\gls{EMG}) sensors' data transmitting via Arduino \gls{BLE} microcontrollers at 50Hz. 

To process this high-frequency data, a custom Android application was engineered to serve as a decentralized edge-computing hub. The software synchronizes the bilateral sensor streams into a 300ms sliding window 
and feeds a dual-tier machine learning pipeline: a cloud-based Random Forest (\gls{RF}) classifier for deep post-session analysis, and a quantized TinyML neural network executing locally for ultra-low latency 
inference.

Empirical evaluation mathematically validated the core sensor fusion hypothesis. Feature importance analysis revealed that localized \gls{EMG} \gls{VAR} and \gls{RMS} outranked standard accelerometer kinematics as 
the primary contributors for differentiating similar basketball movements. Ultimately, the successful deployment of Edge \gls{ML} completely eliminated network-induced latency, enabling instantaneous, active 
injury prevention alerts. This research demonstrates that harmonizing mechanical sensors, electrical physiology, and embedded machine learning provides a good approach to safeguarding athlete health.

%\begin{enumerate}
%  \item What is the problem?
%  \item Why is this problem interesting/challenging?
%  \item What is the proposed approach/solution/contribution?
%  \item What results (implications/consequences) from the solution?
%%\end{enumerate}

% Abstract keywords
\keywords{
  Sports Wearables \and
  Sensor Fusion \and
  Electromyography (\gls{EMG}) \and
  Inertial Measurement Unit (\gls{IMU}) \and
  Edge \gls{ML} \and
  TinyML \and
  Machine Learning (\gls{ML})
} 

